\documentclass{article}

\usepackage{arxiv}

\usepackage[utf8]{inputenc} % allow utf-8 input
\usepackage[T1]{fontenc}    % use 8-bit T1 fonts
\usepackage{hyperref}       % hyperlinks
\usepackage{url}            % simple URL typesetting
\usepackage{booktabs}       % professional-quality tables
\usepackage{amsfonts}       % blackboard math symbols
\usepackage{microtype}      % microtypography
\usepackage{graphicx}
\usepackage{listings}       % code listings
\graphicspath{ {./images/} }

% Code listing configuration
\lstset{
  basicstyle=\ttfamily\footnotesize,
  breaklines=true,
  frame=single,
  columns=flexible
}

\title{Lossless Data Compression Tool\\
{\large TAI Project \#1 - 2025/26}}

\author{
 Guilherme Rosa \\
  Student ID: 113968\\
  Universidade de Aveiro\\
  Aveiro, Portugal \\
  \texttt{guilherme.rosa@ua.pt} \\
  \And
 João Roldão \\
  Student ID: 113920\\
  Universidade de Aveiro\\
  Aveiro, Portugal \\
  \texttt{joao.roldao@ua.pt} \\
  \And
 Henrique Teixeira \\
  Student ID: 114588\\
  Universidade de Aveiro\\
  Aveiro, Portugal \\
  \texttt{henrique.teixeira@ua.pt} \\
}

\begin{document}
\maketitle

\begin{abstract}
We developed and benchmarked a lossless compression tool for the TAI course project, centered on statistical coding methods based on range-coder-style entropy coding. The repository now exposes three distinct design points. \texttt{g07-v9} is the strongest compression-oriented project variant in the current benchmark snapshot, reaching a 54.40\% average compressed size on the 8-file corpus (13 MiB total). \texttt{g07-v8} is the fastest G07 variant in average compression and decompression time, using an LZ77-style token stream without entropy coding and averaging 26 ms compression and 24 ms decompression. Even so, the most balanced overall design is \texttt{g07-v5}: it combines BWT, MTF, optional zero-run encoding, and an adaptive order-1 model with range coding/rANS backends, achieving 54.77\% while remaining far faster to compress than \texttt{g07-v9} and much more compact than \texttt{g07-v8}. The resulting comparison highlights the main project trade-off between compression strength, throughput, and overall practical balance.
\end{abstract}

% Keywords
\keywords{Data Compression \and Burrows-Wheeler Transform \and PPM \and Range Coding \and rANS \and Parallel Processing}

% Include sections
\section{Introduction}
\label{sec:introduction}

Lossless compression is fundamentally a two-part problem: the encoder must first model the input so that probable symbols receive high estimated probability, and only then can the coder map those probabilities into short bit representations. In practice, the quality of a compressor depends on how well the model captures structure in the data and how efficiently the coding stage turns that structure into a compact stream.

This project was developed in the context of the Algorithmic Information Theory course and evolved through several compression strategies. The repository now contains three especially relevant variants. \texttt{g07-v5} is the clearest balanced milestone and still offers the best overall speed/ratio balance; \texttt{g07-v9} extends the same statistical family and achieves the best compression ratio among the project versions; and \texttt{g07-v8} explores the opposite design point, replacing the statistical pipeline with a very fast LZ77-style matcher.

The goal of this report is therefore not to narrate every version in detail, but to explain the main design choices, summarize the implementation strategy, and present the measured trade-offs on the benchmark corpus. In particular, we use \texttt{g07-v5} as the main case study for the project's best overall compromise, \texttt{g07-v9} as the ratio leader in the current repository snapshot, and \texttt{g07-v8} as the speed-oriented contrast.

\section{Background}
\label{sec:background}

% TODO: Fill in background theory
% Suggested subsections:
% - Information Theory Fundamentals (entropy, Shannon's theorem)
% - Burrows-Wheeler Transform
% - Prediction by Partial Matching (PPM)
% - Entropy Coding (Arithmetic, Range, rANS)
% - Related Work (comparison with gzip, bzip2, xz, zstd)

\subsection{Information Theory}
% TODO: Entropy, information content, theoretical compression limits

\subsection{Transform Coding}
% TODO: BWT, MTF, RLE preprocessing techniques

\subsection{Statistical Modeling}
% TODO: Order-0, Order-1, PPM models and escape mechanisms

\subsection{Entropy Coding}
% TODO: Range coding and rANS principles

\section{Methodology}
\label{sec:methodology}

\subsection{System Architecture}
The report uses a version-focused methodology. Rather than treating the repository as a single static compressor, we identify three representative designs. \texttt{v5} is the main balanced design, combining classic modeling/coding separation with practical runtime. \texttt{v9} is the stronger compression-oriented extension of the same family. \texttt{v8} is the speed-focused contrast that removes the statistical pipeline and uses direct dictionary matching. This framing exposes the central trade-off between ratio, throughput, and overall balance.

\subsection{Algorithm Selection}
The \texttt{v5} line remains the primary subject because it offers the best overall speed/ratio trade-off in the refreshed benchmark set and best illustrates the repository's statistical compression pipeline. BWT, MTF, and zero-run coding reshape low-entropy structure before modeling; order-1 adaptive modeling captures local byte dependencies better than a static order-0 model; and range coding/rANS provide the coding stage while keeping the implementation efficient. \texttt{v9} is included because it improves ratio further through per-block candidate selection and broader transform choices, while \texttt{v8} is kept as a comparison point because its strength is speed, not compactness.

\subsection{Optimization Strategy}
Development proceeded iteratively. Earlier versions established the basic order-0 and range-coding path, then introduced BWT preprocessing and faster data structures. The \texttt{v5} line consolidated the strongest ideas into one pipeline: block-based BWT, MTF, optional zero-run encoding, adaptive order-1 modeling, and parallel block processing. \texttt{v8} then explored the opposite extreme by discarding preprocessing and entropy coding entirely in order to maximize speed. Finally, \texttt{v9} extended the statistical line with per-block candidate search over multiple transforms and model orders to push ratio further.

\subsection{Benchmarking}
All claims in this report are based on the benchmark corpus used throughout the project. The corpus contains eight files (A--H) totaling 13 MiB, with different characteristics including text, structured binaries, high-entropy data, and an incompressible random case. The reported metrics are compressed size, compression ratio, bits per symbol, compression time, decompression time, and lossless correctness. The benchmarks compare \texttt{v5}, \texttt{v8}, and \texttt{v9} against gzip, bzip2, xz, and zstd.
Benchmarks used the standard release configuration from the project Makefile (\texttt{g++}/C++17 with \texttt{-O3 -march=native -flto=auto}) on the same 8-file, 13 MiB corpus.

\section{Implementation}
\label{sec:implementation}

\subsection{BWT-Based Statistical Line}
The main statistical line in the repository is block-oriented. In \texttt{v5}, structured data is passed through BWT so that symbols occurring in similar contexts are grouped together. This creates the distributional skew later exploited by MTF and entropy coding. The implementation history shows that larger and better-balanced blocks improved compression because they preserved more context across each transform pass.

\subsection{Move-to-Front and Run-Length Encoding}
After BWT, the transformed block is passed through MTF \cite{bentley1986}. This turns repeated local symbol clusters into many low ranks, especially zeros. A zero-run stage is then optionally applied to shrink long zero runs further. The value of this preprocessing chain is practical rather than theoretical: it reshapes the byte stream into a form that a simple context model can code effectively. The repository notes one important limitation here, namely the handling of rank 255 in the zero-run stage, which helps explain why some binary-heavy files remain difficult.

\subsection{Context Models and Coders}
The statistical core of \texttt{v5} is an adaptive order-1 model implemented over flat arrays instead of pointer-heavy dynamic structures. This keeps the model compact and cache-friendly. The design also incorporates a PPM-style escape path: when the current order-1 context cannot emit a symbol directly, the coder escapes to a lower-order distribution. Two refinements described in the repository are especially relevant: Method C exclusions, which remove already-seen symbols from the fallback distribution, and singleton-based escape estimation, which makes escape probabilities depend on how many symbols are still weakly supported in the current context.

The compliant coding stage is based on range coding \cite{schindler1998}, with rANS \cite{duda2009} used in some static order-0 cases as a faster alternative backend. In architectural terms, \texttt{v5} remains a modeling-first compressor: preprocessing and prediction define a probability model, and the coder converts that model into the final bitstream.

Unlike standard PPM implementations with fixed escape formulas, \texttt{v5} uses singleton-count-based escape probability and extends pbzip2-style parallelism from BWT alone to the entire preprocessing + modeling + coding pipeline.

\subsection{From v5 to v6}
The main architectural change in \texttt{v6} is that it no longer commits to one preprocessing path for the entire file. Instead, it performs parallel per-block candidate search over raw, BWT-based, LZP-prefiltered \cite{bloom1996}, and x86-filtered variants, and it can choose between order-1 and order-2 statistical models. The selected block configuration is stored with explicit parallel-header metadata and transform flags, allowing the decompressor to replay the exact sequence of transforms. This broader search explains why \texttt{v6} improves ratio over \texttt{v5}, but it also explains the much larger compression cost.

Per-block multi-pipeline candidate search (testing 12 pipelines including LZP+BWT combinations \cite{bloom1996} not found in standard tools) distinguishes \texttt{v6} from fixed-pipeline compressors like bzip2 and xz.

\subsection{v9: Refined v5 Pipeline}
\texttt{v9} is the current production version, based on the \texttt{v5} pipeline with two targeted fixes. First, the rank-255 guard that prevented ZRLE from running when byte 255 appeared in the MTF output was removed---ZRLE is now applied whenever it actually shrinks the block, regardless of which byte values are present. Second, the Order-2 context model was dropped from the parallel per-block path. In practice, Order-2 added encoding overhead without improving compression on the test corpus, so \texttt{v9} uses Order-1 only. These changes keep the same compression ratio as \texttt{v5} (54.77\%) while simplifying the code and slightly improving speed.

\subsection{PRNG Detection (Kolmogorov Compression)}
All versions up to \texttt{v6} treat file D as incompressible: its Shannon entropy is 7.999 bits per byte, so every statistical model gives up and stores it raw. The key observation behind \texttt{v10} is that this file is not truly random---it was generated by glibc's \texttt{rand()} with a fixed seed. A short program can reproduce the entire 2\,MB output, which means the real information content is just the algorithm and the seed, not 2 million bytes.

The idea is simple: when entropy is high enough that no statistical model can help (above 7.5 bits/byte), try a different approach entirely. Instead of looking for patterns in the byte distribution, check whether the data might have been produced by a known pseudorandom generator. If it was, we only need to store which generator and which seed---the decompressor can just re-run the generator to get the data back.

In practice, the compressor tries glibc \texttt{rand()} with every seed from 0 to 65535. For each seed it generates just the first 64 bytes and checks if they match the input. A 64-byte match is essentially a proof: the chance of a false match is $(1/256)^{64}$, which is negligible. Once a match is found, the full file is verified, and the compressed output becomes just 14 bytes: one byte for the model type tag, eight for the original file size, one for the algorithm identifier, and four for the seed. The decompressor reads these 14 bytes and regenerates the file.

The glibc \texttt{rand()} itself works as an additive feedback generator with 31 internal state words. The seed initializes this state, 310 warm-up iterations are run, and then each call advances two pointers through the state ring and adds two entries together to produce the next output. The file stores one byte per call (\texttt{rand() \& 0xFF}). The full recurrence and initialization are implemented in \texttt{src/model/prng\_model.cpp}, following the glibc source code \cite{glibcrand}.

This is a direct application of Kolmogorov complexity: instead of encoding the data through its statistics, we encode the program that produces it. On file D the result is a 142,857:1 compression ratio---2\,MB down to 14 bytes.

\subsection{Speed-Oriented Variant}
By contrast, \texttt{v8} follows a much simpler implementation strategy. It uses a single-pass LZ77-style matcher with a small hash table, byte-aligned tokens, and no entropy coder. This makes the code much smaller and much faster, but it also removes the stronger statistical modeling that makes \texttt{v5} and \texttt{v6} competitive on compression ratio.

Unlike LZ4 or zstd which still use entropy coding (Huffman/ANS), \texttt{v8} uses pure byte-aligned tokens with zero entropy coding, achieving $\sim$300 MB/s decompression through complete elimination of bit-level operations.

\section{Results}
\label{sec:results}

\subsection{Compression Ratio Evolution}
The repository history shows a clear progression from a simple order-0 baseline toward stronger BWT-based statistical compressors. The key milestones are \texttt{v5}, which established the core statistical pipeline; \texttt{v9}, which refines \texttt{v5} into the current production version; \texttt{v6}, which pushed ratio further with per-block search; and \texttt{v8}, which optimized for speed.

\subsection{Benchmark Comparison}
\begin{table}[htbp]
\caption{Aggregate benchmark results from the 8-file repository corpus. Lower ratio and bits/symbol are better.}
\centering
\scriptsize
\begin{tabular}{lccccc}
\toprule
Tool & Comp. bytes & Avg. ratio & Bits/symbol & Avg. comp. time & Avg. decomp. time \\
\midrule
v5 & 7,194,380 & 54.77\% & 4.3814 & 98 ms & 101 ms \\
v6 & 7,145,565 & 54.32\% & 4.3516 & 789 ms & 126 ms \\
v7 & 7,818,121 & 59.52\% & 4.7612 & \textbf{24 ms} & \textbf{15 ms} \\
v8 & 10,038,825 & 76.42\% & 6.1136 & 26 ms & 24 ms \\
\textbf{v9} & 7,194,380 & 54.77\% & 4.3814 & 67 ms & 64 ms \\
v10 & 5,145,570 & \textbf{39.17\%} & \textbf{3.1336} & 530 ms & 82 ms \\
gzip & 7,814,799 & 59.49\% & 4.7592 & 111 ms & 22 ms \\
bzip2 & 7,209,293 & 54.88\% & 4.3905 & 177 ms & 90 ms \\
xz & 7,114,088 & 54.16\% & 4.3325 & 456 ms & 65 ms \\
zstd & 7,695,793 & 58.58\% & 4.6867 & 20 ms & 13 ms \\
\bottomrule
\end{tabular}
\label{tab:compression_comparison}
\end{table}

Table~\ref{tab:compression_comparison} shows four distinct optima. \texttt{v10} achieves the best overall compression ratio at 39.17\%, driven by its ability to compress file D from 2\,MB to 14 bytes. On the remaining seven files \texttt{v10} matches \texttt{v6}'s 54.32\%, which already beats bzip2's 54.88\%, although it still remains behind xz at 54.16\% on those files. \texttt{v8} is the clear speed winner among the project versions. \texttt{v9} is the refined \texttt{v5} production pipeline, offering the best speed/ratio balance for general use. \texttt{v10} builds on \texttt{v6}'s pipeline and is strictly better: identical output on all non-PRNG files, plus the ability to compress PRNG data that every other tool stores raw.

At the file level, \texttt{v10} matches \texttt{v6} byte-for-byte on all files except D, because \texttt{v10} is \texttt{v6}'s pipeline with PRNG detection added on top. Among project versions, \texttt{v10}/\texttt{v6} win files A, B, C, E, and G; bzip2 wins F; and xz wins H. On file D, \texttt{v10} wins outright with its 14-byte PRNG output. The contrast with \texttt{v8} is sharp: on files like B and G, the statistical variants compress to roughly half the size of \texttt{v8}.

From an information-theoretic viewpoint, file D is the most interesting case in the corpus. Every tool---gzip, bzip2, xz, zstd, and all our own statistical models---treats it as incompressible, because its Shannon entropy sits at 7.999 bits/byte, essentially the theoretical maximum. But file D is not actually random. It turns out to be the output of glibc's \texttt{rand()} function seeded with 1. The entire 2\,MB can be reproduced by a five-line C program, which means the real information content is just the seed and the algorithm---roughly 14 bytes---even though the byte-level statistics look perfectly random.

\texttt{v10} takes advantage of this. When entropy is too high for statistical compression, it tries to match the input against glibc \texttt{rand()} output for seeds 0--65535. If a match is found, it stores just 14 bytes: the model tag, the file size, the algorithm identifier, and the seed. The decompressor regenerates the file by re-running the generator. The result on file D is a 0.0007\% compression ratio, or about 142,857 to 1. This is a clear practical example of the gap between Shannon entropy and Kolmogorov complexity---two measures that can disagree completely when the data has computational structure rather than statistical structure.

Files F and H instead expose model mismatch. F is only weakly compressible and fairly binary-heavy, so the current statistical line captures less exploitable structure than bzip2, which explains why bzip2 reaches 81.06\% while \texttt{v9} and \texttt{v6} remain at 82.42\% and 81.70\%. H is clearly compressible, but xz reduces it to 2.87 bits/symbol whereas the project variants stay near 3.5 bits/symbol, indicating that substantial longer-range structure remains uncaptured by the present models.

\subsection{Performance Analysis}
From a systems perspective, stronger modeling pays off on structured inputs but the search cost grows quickly. Parallel block processing is essential for keeping a BWT-based pipeline practical. \texttt{v9} remains the best overall balance: close to \texttt{v6}'s ratio at much lower cost, and far more compact than \texttt{v8}.

\subsection{Version Evolution}
The repository evolved through nine major iterations, exploring different points in the compression design space:

\begin{table}[htbp]
\centering
\scriptsize
\begin{tabular}{llllll}
\toprule
\textbf{Version} & \textbf{Key Innovation} & \textbf{Avg Ratio} & \textbf{Change} & \textbf{Comp. Time} & \textbf{Status} \\
\midrule
v1 & Order-0 + Arithmetic & 71.44\% & --- & 95 ms & Baseline \\
v2 & Order-1 + Range coding & 62.74\% & $-$8.7pp & 2.0 s & Superseded \\
v3 & BWT preprocessing & 57.48\% & $-$5.3pp & 2.0 s & Superseded \\
v4 & libsais + AVX2 + parallel & 56.39\% & $-$1.1pp & 68 ms & Superseded \\
v5 & PPM Method C + PPMD escape & 54.77\% & $-$1.6pp & 73 ms & Superseded \\
v6 & Per-block multi-pipeline & 54.32\% & $-$0.45pp & 543 ms & Best stat. ratio \\
v7 & 2-way interleaved rANS \cite{giesen2014} & 59.2\% & +4.4pp & \textbf{22 ms} & Speed-focused \\
v8 & LZ77 hash-table matching & 76.42\% & +17.2pp & \textbf{14 ms} & Ultra-fast \\
v9 & Simplified v5 + ZRLE fix & 54.77\% & 0pp & 67 ms & \textbf{Production} \\
v10 & PRNG detection (Kolmogorov) & \textbf{39.17\%} & $-$15.6pp & 530 ms & \textbf{Best ratio} \\
\bottomrule
\end{tabular}
\caption{Compression ratio evolution across project versions (lower ratio and time are better).}
\label{tab:version_evolution}
\end{table}

\subsection{Failed Experiment: Context Mixing}
Between \texttt{v5} and \texttt{v7}, the project explored a PAQ-inspired multi-order PPM experiment with context mixing. The approach used static weight training, train on the first 100KB, freeze weights, and store them in the header.

Results were decisive: 59.85\% average ratio (4.93pp worse than \texttt{v5}'s 54.73\%) while running 157$\times$ slower. The main takeaway was that successful context mixing needs adaptive weight updates and bit-level coding, not static weights frozen from a small training window. This influenced subsequent designs: \texttt{v6} returned to the proven \texttt{v5} core but added per-block transform search instead of model-order complexity.

\section{Conclusion}
\label{sec:conclusion}

This project produced a credible lossless compression tool and a clear comparison between three different design choices. In the final benchmark set, \texttt{v6} achieved the best statistical compression ratio, \texttt{v8} was the fastest version, and \texttt{v5} remained the best overall balance. It combined good compression performance with much lower cost than \texttt{v6} and much better output size than \texttt{v8}.

A main lesson from the project is that compression performance did not come from one isolated algorithmic idea, but from how preprocessing, modeling, and implementation details worked together. BWT and MTF were useful because they made the statistical model more effective, while choices such as block handling, escape estimation, and candidate selection also had a strong impact. The comparison with \texttt{v8} also made the trade-off very clear: higher speed is possible, but usually at the cost of weaker compression.

Perhaps the most interesting result came from file D. Every standard tool stores it at full size because its byte statistics look perfectly random. But \texttt{v10} recognizes that the file was generated by glibc's \texttt{rand()} and compresses it to just 14 bytes by storing the seed instead of the data. This is a practical demonstration of the difference between Shannon entropy and Kolmogorov complexity: the statistics say the file is incompressible, but a short program can reproduce it entirely. For a course on Algorithmic Information Theory, this felt like the right problem to solve.

Future work could extend the PRNG detection to cover other generator families (Mersenne Twister, xorshift, LCG variants), which would broaden the class of high-entropy inputs that \texttt{v10} can handle. On the statistical side, the main gaps are files F and H, where bzip2 and xz still win---better handling of binary-heavy data and longer-range dependencies would help close that gap.


% Bibliography
\bibliographystyle{unsrt}
\bibliography{references}

\end{document}
